\section{Orchestration and Obligation Liveness}
\label{sec:sec14}\label{sec:temporal}

The framework fixes \emph{what} must hold---conservation, atomicity, idempotency, purity---and leaves open \emph{how} lifecycle events are scheduled, retried, and sequenced. Safety alone does not move the system forward: a coupon due on a date causes no error if it never fires. This section states the requirements an execution engine must meet to close that gap, maps the executor onto it, and then develops obligation liveness as a first-class subsystem with its own invariants (P21--P23) and proof.

%----------------------------------------------------------------------
\subsection{Execution-engine requirements}
\label{sec:exec-requirements}
%----------------------------------------------------------------------

Four requirements follow from the framework's own primitives. They bind any execution engine; the engine is an implementation choice that satisfies them, not a result derived here.

\begin{enumerate}[nosep]
\item \textbf{At-least-once delivery.} Idempotency (P5, transaction level; P6, lifecycle level) makes duplicate execution harmless but does not make an event fire. The engine must dispatch every scheduled event and re-dispatch on worker failure. At-least-once delivery composed with idempotent activities yields exactly-once semantics.

\item \textbf{Durable timers surviving restarts.} A coupon scheduled for a fixed future date fires on that date irrespective of process restarts, deployments, or infrastructure failures in the interval. Timer state is persisted by the engine, not held in process memory.

\item \textbf{Deterministic replay.} Time travel (Property~6) requires ledger state at any past instant to be reconstructable by replaying the move stream. The engine carries the complementary property for orchestration: workflow state at any point is reconstructable by replaying an append-only workflow history. Both derive correctness from one principle---deterministic replay of an immutable log.

\item \textbf{Single-threaded execution per instance.} Each workflow instance processes signals, timer fires, and activity completions one at a time, in deterministic order. This is the single-writer model; it eliminates concurrency races without a command queue or consensus protocol.
\end{enumerate}

\noindent The reference architecture adopts Temporal.io, whose durable-execution model satisfies all four. Thereafter ``the engine'' denotes the adopted instance, and its guarantees enter the liveness proof as named lemmas (\S\ref{sec:liveness-proof}).

%----------------------------------------------------------------------
\subsection{The executor as an activity}
\label{sec:executor-activity}
%----------------------------------------------------------------------

The executor---the single component permitted to mutate the ledger---validates conservation, checks idempotency by transaction ID (P5), and commits atomically. It runs as a retriable engine \emph{activity}: a unit of work the engine re-invokes on transient failure.

The mapping is exact on two points. \textbf{Retry.} The engine retries on transient failures (network partition, process crash, transient database error); the executor's transaction-ID idempotency check makes a retried commit a no-op, upgrading at-least-once dispatch to exactly-once effect. Conservation violations, referential-integrity errors, and idempotency rejections are non-retryable---they signal logic errors, never transient ones, and a retry never succeeds. \textbf{Idempotency contract.} The activity's idempotency key is the transaction ID. If that ID is already committed, the executor returns the prior result without re-applying moves, satisfying the engine's requirement that an activity called repeatedly with the same input produce one observable effect.

%----------------------------------------------------------------------
\subsection{The due-event scheduler}
\label{sec:temporal-scheduler}
%----------------------------------------------------------------------

Every contractual obligation with a deterministic future date is mapped, at unit registration, to a durable timer: bond coupons, futures daily settlement, option expiry, IRS resets, CSA margin calls, managed-account resets (Table~\ref{tab:obligation-types}, deterministic-date rows). The timer is persisted by the engine and survives server restarts, worker deployments, version upgrades, and data-centre failovers. This resolves the liveness gap for dated obligations: the timer fires regardless of infrastructure failures, and idempotency guarantees it fires at most once.

A due event does not proceed on stale or unavailable market data. Each lifecycle workflow runs a data-quality check before invoking the lifecycle function; on failure it defers the event and retries after a fixed interval, the deferral itself durable across worker crashes. Market-data staleness thresholds and cross-source validation are specified in the market-data satellite specification; this section consumes that gate, it does not define it.

%----------------------------------------------------------------------
\subsection{Concurrency: single-writer by construction}
\label{sec:temporal-concurrency}
%----------------------------------------------------------------------

One lifecycle workflow is assigned per unit (or per wallet for per-wallet state such as futures accumulated cost). Because each instance processes events single-threaded in deterministic order, concurrent lifecycle events on the same unit serialise: two margin calls on one CSA cannot interleave, two coupons on one bond cannot race, and a daily settlement and an intraday trade on one futures contract process in strict causal order. The executor therefore never receives concurrent mutations for a unit from the orchestration layer; its idempotency check (P5) is a second line of defence against a duplicate reaching it through an activity retry.

Operations spanning units---corporate actions, close-out netting, CSA-level margin---are coordinated by a parent workflow whose deterministic execution fixes causal order across children; no shared mutable state exists between them. Serialisation is per unit, not global: distinct units execute in parallel, so throughput scales with the number of active units and the worker pool, not with any single unit's lifecycle complexity.

%----------------------------------------------------------------------
\subsection{Obligation liveness}
\label{sec:obligation-liveness}
%----------------------------------------------------------------------

The scheduler closes the gap only for obligations whose deadline is known at registration. A second class---\emph{event-triggered} obligations---arises during the lifecycle with no timer at registration: an SBL recall creates a return-by deadline absent at loan inception; a collateral-eligibility change creates a substitution demand with a cut-off; a margin call creates a delivery obligation at T+1 from the call; a counterparty default triggers close-out netting on master-agreement deadlines. Handled only inside workflow code, such obligations are invisible to the framework, which cannot then distinguish ``tracked'' from ``forgotten by the workflow author.'' The obligation as a first-class object closes this gap.

\subsubsection{The obligation type}
\label{sec:obligation-type}

\begin{definition}[Obligation]
\label{def:obligation}
An \emph{obligation} is a tuple $o = (\mathit{id},\, \mathit{type},\, \mathit{source},\, t_d,\, D,\, \kappa)$:
\begin{itemize}[nosep]
\item $\mathit{id}$: a unique identifier, derived deterministically from the source event and obligation type.
\item $\mathit{type}$: a member of the obligation-type enumeration (Table~\ref{tab:obligation-types}).
\item $\mathit{source}$: the unit, agreement, or regulatory event that created it.
\item $t_d$: the \emph{deadline} by which the obligation must be discharged.
\item $D$: the \emph{discharge predicate} on ledger state. The obligation is discharged when a move set $M$ journaled before $t_d$ gives $D(\mathrm{state\_after}(M)) = \mathit{true}$.
\item $\kappa$: the \emph{compensation action}, a deterministic function computing compensating moves when the deadline passes undischarged.
\end{itemize}
\end{definition}

\noindent The obligation record and its discharge predicate are expressed directly in the type system. The discharge predicate $D$ and the compensation $\kappa$ are \emph{total} functions of ledger state---same state, same verdict---so a discharge decision is deterministic and replayable. $\kappa$ returns \texttt{Just}~moves when it can compute valid compensating moves and \texttt{Nothing} when it cannot, the latter forcing \textsc{Defaulted}.

\begin{lstlisting}[language=Haskell]
-- D and kappa are total functions of ledger state (deterministic verdicts).
data Obligation = Obligation
  { obId         :: ObId                        -- derived from source event + type
  , obType       :: ObType                      -- taxonomy (obligation-types table)
  , obSource     :: Source                      -- unit / agreement / regulatory event
  , obDeadline   :: Time                        -- t_d
  , obDischarge  :: LedgerState -> Bool         -- D : the discharge predicate
  , obCompensate :: LedgerState -> Maybe [Move] -- kappa : Just on success, Nothing fails
  }
\end{lstlisting}

\noindent The lifecycle state is split so that the type, not a convention, forbids leaving a terminal state: \textsc{Pending} and \textsc{Attempted} are the only \emph{live} states, and the handler consumes a live state alone. \textsc{Discharged}, \textsc{Compensated}, and \textsc{Defaulted} are terminal and have no exit edge---``an obligation leaves a terminal state'' is unrepresentable, not merely unreached.

\begin{lstlisting}[language=Haskell]
data Live     = Pending | Attempted                  -- non-terminal
data Terminal = Discharged | Compensated | Defaulted  -- terminal: no exit edge
data ObState  = Active Live | Settled Terminal        -- stored per obligation

data Trigger  = DischargeSignal | DeadlineFired       -- the two things that move o

-- Total over (Live x Trigger). `step` takes only Live, so a Terminal can never
-- be transitioned. On DeadlineFired every arm returns Right (a terminal state):
-- that absence of a Left IS lemma L5 (handler totality), the handler half of P21.
step :: Obligation -> LedgerState -> Live -> Trigger -> Either Live Terminal
step o s _ DischargeSignal
  | obDischarge o s = Right Discharged
  | otherwise       = Left  Attempted
step o s _ DeadlineFired
  | obDischarge o s = Right Discharged
  | otherwise       = case obCompensate o s of
      Just _ms      -> Right Compensated    -- kappa succeeded
      Nothing       -> Right Defaulted      -- kappa failed
\end{lstlisting}

\noindent No categorical name is introduced here: a record, a sum type split into live and terminal, and one total function discharge the obligations the type must meet; nothing further would buy a removed bug or removed code. The state-table transitions of \S\ref{sec:obligation-type} are exactly the cases of \texttt{step}. Behaviour, made concrete on the CSA call below (\S\ref{sec:obligation-example-csa}): \texttt{step csaVM} on a post-delivery state with \texttt{DischargeSignal} returns \texttt{Right Discharged}; on a no-delivery state with \texttt{DeadlineFired} it returns \texttt{Right Compensated}.

\begin{center}
\small
\begin{tabular}{l l l}
\toprule
\textbf{From} & \textbf{To} & \textbf{Condition} \\
\midrule
\textsc{Pending}   & \textsc{Attempted}   & discharge attempted \\
\textsc{Pending}   & \textsc{Discharged}  & $D = \mathit{true}$ (early) \\
\textsc{Pending}   & \textsc{Compensated} & $t > t_d$, $\kappa$ fires \\
\textsc{Attempted} & \textsc{Discharged}  & $D = \mathit{true}$ \\
\textsc{Attempted} & \textsc{Compensated} & $t > t_d$, $\kappa$ fires \\
\textsc{Pending}   & \textsc{Defaulted}   & $t > t_d$, $\kappa$ fails \\
\textsc{Attempted} & \textsc{Defaulted}   & $t > t_d$, $\kappa$ fails \\
\bottomrule
\end{tabular}
\end{center}

\subsubsection{Taxonomy}
\label{sec:obligation-taxonomy}

\begin{table}[ht]
\centering
\small
\begin{tabular}{l l l l}
\toprule
\textbf{Obligation type} & \textbf{Scope} & \textbf{Trigger} & \textbf{Compensation $\kappa$} \\
\midrule
\multicolumn{4}{l}{\textit{Deterministic-date (timer at registration)}} \\
Bond coupon & Unit & Timer at coupon date & Default event \\
Option expiry & Unit & Timer at termination date & N/A (auto-expire) \\
IRS reset & Unit & Timer at reset date & Default event \\
Futures daily VM & Unit & Cron at exchange EOD & Position liquidation \\
SBL term loan return & Unit & Timer at maturity & Buy-in (GMSLA~9.3) \\
\midrule
\multicolumn{4}{l}{\textit{Event-triggered (timer at event time)}} \\
SBL recall return & Unit & Recall signal & Buy-in (GMSLA~9.3) \\
SBL manuf.\ dividend & Unit & Corp.\ action record date & Cost attribution \\
SBL collateral subst. & Unit & Demand / eligibility change & Shortfall escalation \\
CSA VM delivery & Agreement & Margin call computation & Close-out netting (ISDA) \\
CSA IM delivery & Agreement & Regulatory schedule & Close-out netting \\
Collateral substitution & Agreement & Schedule amendment & Close-out netting \\
SBL collateral top-up & Unit & MTM threshold breach & Recall / default \\
Close-out netting & Agreement & Default event & Waterfall / loss alloc. \\
\midrule
\multicolumn{4}{l}{\textit{Regulatory}} \\
SFTR report & Unit & SBL lifecycle event & Regulatory penalty \\
SLATE report & Unit & SBL lifecycle event & Regulatory penalty \\
EMIR report & Unit & Trade/lifecycle event & Regulatory penalty \\
Settlement instruction & Unit & Settlement-type transaction & Failed settlement \\
\bottomrule
\end{tabular}
\caption{Obligation types by trigger and scope. Deterministic-date obligations are scheduled at unit registration; event-triggered obligations are created when the triggering event is processed; regulatory obligations are created alongside the business event.}
\label{tab:obligation-types}
\end{table}

\subsubsection{The obligation store}
\label{sec:obligation-store}

\begin{definition}[Obligation store]
\label{def:obligation-store}
The \emph{obligation store} $\mathcal{O}$ is the set of obligations recovered from the event log by filtering to obligation-typed entries. Each obligation is added atomically within the transaction that creates its triggering event; each state transition is a logged state-only transaction.
\end{definition}

\noindent The obligation store is a projection of the event log, not a second source of truth. Each change to an obligation's state is a logged transaction (Definition~\ref{def:obligation-store}); the stored value is a read cache, overwritten in place. Replaying the log to any instant rebuilds the value that held then, so overwriting loses nothing, and a logged transaction is the only thing that changes the value. The value exists from the obligation's creation onward.

\subsubsection{Registration}
\label{sec:obligation-registration}

Obligations register through two channels. At \textbf{unit registration}, the lifecycle function extracts dated obligations from the CDM \texttt{EconomicTerms} (coupon, reset, termination schedules), one per date; these correspond one-to-one with the scheduler's timers (\S\ref{sec:temporal-scheduler}), making explicit what the timer alone left implicit---each has a defined discharge predicate and compensation action. At \textbf{event processing}, a lifecycle event that creates an obligation returns it alongside its moves and state deltas; the executor commits the obligation registration atomically with the triggering event.

\begin{principle}[Obligation completeness]
\label{princ:obligation-completeness}
A lifecycle function processing an obligation-creating event includes the obligation in its output. A function that omits a required obligation is incorrect, irrespective of whether its moves are valid.
\end{principle}

\noindent The CDM event-type enumeration bounds the test space, so this principle is enforceable by property-based testing: for each obligation-creating event type (Table~\ref{tab:obligation-types}), the generator produces events and the postcondition verifies a non-empty, well-formed obligation list. A typed return that mandates an obligation list for obligation-creating event types makes omission a compile-time error.

\subsubsection{The obligation workflow}
\label{sec:obligation-workflow}

Each obligation is managed by one workflow---spawned per obligation, or embedded in the unit or agreement lifecycle workflow that owns its source. The pattern is uniform: a selector waits on a discharge signal channel and a deadline timer. On a discharge signal, the handler checks $D$ against ledger state; if satisfied, the obligation becomes \textsc{Discharged} and the workflow completes, otherwise it becomes \textsc{Attempted} and waits. When the timer fires at $t_d$, the handler re-checks $D$; if satisfied, \textsc{Discharged}; otherwise it commits $\kappa$ through the executor, reaching \textsc{Compensated} on success and \textsc{Defaulted} (raising a default event) on failure. The workflow ID is derived deterministically from $o.\mathit{id}$, so the engine's duplicate-start rejection prevents double-spawning.

\subsubsection{Liveness invariants}
\label{sec:obligation-invariants}

\begin{invariant}[P21---Obligation Liveness]
\label{inv:P21}\label{inv:obligation-liveness}
For every $o \in \mathcal{O}$ with deadline $t_d$:
\[
\forall\, t > t_d : \quad o.\mathit{state} \in \{\textsc{Discharged},\; \textsc{Compensated},\; \textsc{Defaulted}\}.
\]
No obligation persists in \textsc{Pending} or \textsc{Attempted} beyond its deadline.
\end{invariant}

\begin{invariant}[P22---Obligation Conservation]
\label{inv:P22}\label{inv:obligation-conservation}
Every discharging move set and every compensation move set for $o$ satisfies the conservation law $Q(u) = 0$ (P1). The obligation mechanism does not bypass the executor.
\end{invariant}

\begin{invariant}[P23---Obligation Idempotency]
\label{inv:P23}\label{inv:obligation-idempotency}
Discharging the same obligation twice, by $o.\mathit{id}$, has no incremental effect: $D$ is checked before moves are emitted, and if already satisfied no moves are produced. This composes with transaction-ID idempotency (P5) and lifecycle idempotency (P6).
\end{invariant}

\subsubsection{Proof of liveness}
\label{sec:liveness-proof}

\begin{theorem}[Obligation liveness]
\label{thm:obligation-liveness}
Under the executor, engine, and obligation store, P21 (Invariant~\ref{inv:P21}) holds: no obligation persists in \textsc{Pending} or \textsc{Attempted} beyond its deadline.
\end{theorem}

\begin{proof}
Five lemmas compose.

\textbf{L1 (Registration completeness).} Every obligation-creating event registers an obligation in $\mathcal{O}$. By Principle~\ref{princ:obligation-completeness}, registration is part of the lifecycle function's output contract; the executor commits it atomically with the triggering event; property-based testing verifies a non-empty obligation list per obligation-creating event type.

\textbf{L2 (Workflow creation).} Every $o \in \mathcal{O}$ is managed by a workflow with a timer at $t_d$. The executor's post-commit hook spawns it; the deterministic workflow ID makes the engine reject double-spawning.

\textbf{L3 (Timer durability).} The timer at $t_d$ is persisted by the engine and survives server restarts, worker deployments, version upgrades, and data-centre failovers---a guarantee of the engine (\S\ref{sec:exec-requirements}).

\textbf{L4 (Timer fires).} The timer fires at $t_d$, by the engine's durable-execution guarantee. \textit{Prerequisite}: the engine cluster is available at $t_d$. Cluster availability is an operational prerequisite, not a framework property; multi-region replication mitigates but does not eliminate the dependency.

\textbf{L5 (Handler totality).} On firing, the handler transitions $o$ out of \textsc{Pending}. The handler is total: $D = \mathit{true}$ yields \textsc{Discharged}; otherwise $\kappa$ fires, yielding \textsc{Compensated} on success and \textsc{Defaulted} (with a default event) on failure. Every branch terminates in a terminal state.

\textbf{Composition.} For $o$ with deadline $t_d$: $o \in \mathcal{O}$ (L1); a workflow with a timer at $t_d$ exists (L2); the timer is durable (L3) and fires (L4); the handler reaches a terminal state (L5). Hence for all $t > t_d$, $o.\mathit{state} \in \{\textsc{Discharged}, \textsc{Compensated}, \textsc{Defaulted}\}$.
\end{proof}

\noindent P22 follows from P1 and executor exclusivity: every discharging and compensating move passes through the executor, which validates $Q(u)=0$ before commit. P23 composes with P5 and P6: the obligation ID is part of the transaction source metadata, the discharge predicate prevents re-emission, and the transaction ID prevents re-commitment. The composition is summarised below.

\begin{center}
\small
\begin{tabular}{l l}
\toprule
\textbf{Property} & \textbf{Guarantee source} \\
\midrule
Every obligation fires & engine durable timer (L4) \\
No obligation fires twice & idempotency (P5, P6, P23) \\
Discharge preserves conservation & executor (P1, P22) \\
Handler reaches terminal state & workflow totality (L5) \\
\bottomrule
\end{tabular}
\end{center}

\noindent Every obligation---including those created dynamically during the lifecycle---is thereby tracked, timed, and resolved, not only those with a deterministic registration-time date.

\subsubsection{Worked example: CSA variation margin call}
\label{sec:obligation-example-csa}

Alice and Bob hold a CSA over five IRS trades, threshold \$500{,}000, minimum transfer amount (MTA) \$250{,}000. The daily margin workflow computes aggregate mark-to-market of \$2.3M in Alice's favour; required margin is $\max(0, |\$2.3\text{M}| - \$500\text{K}) = \$1.8\text{M} \geq \$250\text{K}$, so a call is generated. The call registers an obligation, committed atomically with the call event:

\begin{lstlisting}
Obligation(
    id:            "CSA-AB-20260328-VM",
    type:          CSA_VARIATION_MARGIN,
    source:        CSA-AB,
    deadline:      2026-03-31T17:00 (T+1, next business day COB),
    discharge:     Bob.coll_post(USD, CSA-AB) increased by >= 1,800,000,
    compensation:  CSA_CloseOutNetting(CSA-AB)
)
\end{lstlisting}

\noindent \textbf{Discharge.} Bob delivers \$1.8M before the deadline as a conservation-preserving move; the discharge signal satisfies $D$; obligation $\to$ \textsc{Discharged}. \textbf{Compensation.} Bob does not deliver; the timer fires; $D$ fails; \texttt{CSA\_CloseOutNetting} terminates all five trades, computes close-out values per the ISDA Master Agreement, and emits one net settlement move through the executor; obligation $\to$ \textsc{Compensated}. Conservation on the discharge path:
\begin{align*}
\Delta\, w_{\mathrm{Bob}}(\text{USD})[\mathrm{own}] &= -1{,}800{,}000, &
\Delta\, w_{\mathrm{Alice}}(\text{USD})[\mathrm{coll\_recv}] &= +1{,}800{,}000, \\
Q(\text{USD}) &= 0. \checkmark
\end{align*}

\subsubsection{Worked example: SBL collateral substitution}
\label{sec:obligation-example-sbl}

A lender demands substitution: bond XYZ is no longer eligible under the updated collateral schedule. The eligibility check registers an obligation (IBP references denote clauses of the ISLA Best Practice Handbook~\cite{isla-ibp}):

\begin{lstlisting}
Obligation(
    id:            "SBL-L123-SUBST-20260328",
    type:          SBL_COLLATERAL_SUBSTITUTION,
    source:        loan L123,
    deadline:      2026-03-28T15:30 (1h before DVP cut-off, IBP-170),
    discharge:     new eligible collateral posted AND old released,
    compensation:  CollateralShortfallEscalation(L123)
)
\end{lstlisting}

\noindent \textbf{Discharge.} The borrower delivers German Bund as replacement in a four-move transaction, each move satisfying the Single-Coordinate Move Principle (Principle~\ref{princ:single-coord}):

\begin{lstlisting}
tau_subst = Transaction(type = SETTLEMENT):    -- txIntroduce = Nothing; applyTx is the sole door
  Move 1: Borrower coll_post(XYZ)  -= old_amt   -- release old
  Move 2: Lender   coll_recv(XYZ)  -= old_amt
  Move 3: Borrower coll_post(BUND) += new_amt   -- post new
  Move 4: Lender   coll_recv(BUND) += new_amt
\end{lstlisting}

\noindent Per-unit conservation: $\Delta Q(\text{XYZ}) = -\mathit{old\_amt} + \mathit{old\_amt} = 0$; $\Delta Q(\text{BUND}) = \mathit{new\_amt} - \mathit{new\_amt} = 0$. Obligation $\to$ \textsc{Discharged}. \textbf{Compensation.} The cut-off passes without delivery; the timer fires; \texttt{CollateralShortfallEscalation} marks L123 collateral-deficient in the SBL unit state, raises an exposure-breach notification, and---if the breach persists beyond the GMSLA grace period---triggers recall or default per GMSLA~\S10. Obligation $\to$ \textsc{Compensated}.

\subsubsection{Relationship to SBL invariants}
\label{sec:obligation-sbl-relationship}

P21 supplies the temporal guarantee underpinning the SBL invariants (P11--P20). It guarantees collateral top-up obligations resolve by deadline---delivery or compensation---so P13 (Collateral Sufficiency) cannot be violated by an untracked call. It guarantees each reportable event's reporting obligation fires before its T+1 deadline or escalates, sustaining P16 (SFTR Completeness). It guarantees recall obligations resolve, so shares cannot remain on loan past a recall deadline without return or buy-in, sustaining P18 (Lender Ownership Invariance).

\subsubsection{Property-based testing for obligations}
\label{sec:obligation-testing}

Three properties extend the testing framework. \textbf{Registration completeness}: for random obligation-creating events, the output includes a non-empty obligation list with valid deadline, discharge predicate, and compensation action. \textbf{Discharge conservation}: every discharge move set satisfies $Q(u) = 0$. \textbf{Handler totality}: over random obligation states $\times$ \{discharge-before-deadline, no-discharge, compensation-success, compensation-failure\}, the handler reaches a terminal state in every case. The CDM event-type enumeration bounds the obligation-creating event space, making coverage structurally complete.

\subsubsection{Assumptions and limitations}
\label{sec:obligation-limitations}

Theorem~\ref{thm:obligation-liveness} depends on four assumptions, each at the framework boundary. \textbf{Engine availability at $t_d$}: an unavailable cluster delays the timer until recovery; liveness holds eventually, not instantaneously. \textbf{Executor availability}: an unavailable executor leaves the obligation \textsc{Pending} under retry until recovery---a bounded delay, not a permanent failure. \textbf{External cooperation}: discharge conditions requiring counterparty action cannot be compelled; compensation is the contractual fallback. \textbf{Lifecycle-function correctness}: L1 assumes lifecycle functions register all required obligations, enforceable by property-based testing and, in a typed implementation, by a mandatory obligation list in the return type. These match the framework's existing boundary: correctness within the ledger, given infrastructure availability and correct lifecycle functions; external actors and infrastructure availability lie outside it.
